\chapter{Urbanismo industrial}
\section{Planificación industrial}

\section{Localización}

\section{Proyecto}

\subsection{Implantación}
\subsubsection{Polígonos industriales}
``Parque Industrial'' o ``Polígono Industrial'': Agrupación de varias industrias con el fin de obtener las ventajas que puede suponer el estar dotadas de una serie de servicios comunes.

Son agrupaciones idóneas para el emplazamiento de plantas industriales correspondientes a pequeñas y medianas empresas, ya que entre los servicios comunes pueden tener no solo los correspondientesd a la infraestructura técnica o de equipamiento (abastecimiento de energía, agua, estaciones depuradores, desagües, etc.), sino que incluso pueden disponer de servicios empresariales comunes que pueden mejorar la eficacia individual de cada fábrica de las integradas en el Polígono.

\subsubsection{Parques científico tecnológicos}
Nuevos espsacios tecnológicos basados en empresas de alta tecnología. Se pueden crear en lugares remotos o de baja industrialización cuando se den unos factores decisivos:
\begin{itemize}
    \item Presencia de instituciones dedicadas a la investigación y a la formación
    \item Incentivos fiscales y financieros
    \item Disponibilidad de terreno industrial
    \item Mercado de trabajo local con ingenieros y técnicos de calidad
    \item Buenos sistemas de transportes y comunicaciones
    \item Calidad del entorno, flexibilidad burocrática y buena imagen de la ubicación
\end{itemize}

Suelen estar compuestos de 3 patas: Los centros de investigación y universidades, grandes empresas y PYMEs.

\subsubsection{Parámetors urbanísticos}
Las normas de planeamiento estructurantes o de desarrollo, según el caso, establecerán distintas ordenanzas fijando una serie de datos o parámetros urbanísticos asociados a tipologías edificatoria, usos, ubicación, etc..., y que deberá cumplir la edificación que se pretenda ser realizada en una parecela concreta:
\begin{itemize}
    \item Ocupación en planta.
    Concepto: \% en planta de la parcela que se puede ocupar, como máximo, con edificaciones.
    Limitaciones: se suelen establecer máximos, que no tienen que ser iguales para cada parcela.
    Dependiendo de otras restricciones diferentes a ésta (retranqueos), puede que no lleguemos a poder aprovechar este máximo en algunas ocasiones.
    Este parámetro no indica cuál es el máximo total (en diferentes plantas) que se puede edificar.
    \item Retranqueo.
    Concepto: son las distancias que deben dejar libres entre los diversos linderos de la parcela y las edificaciones.
    Limitaciones: las ordenanzas suelen establecer unos retranqueos mínimos que hay que respetar.
    \item Edificabilidad.
    Concepto: superficie máxima que se puede construir en las diversas plantas de los diferentes edificios que se incluirán en la parcela.
    Limitaciones: normalmente se define de forma directa, como máximo de $m^2$ construidos por cada $m^2$ de parcela ($m^2 / m^2$).
    Dependiendo de otras restricciones (altura máxima), puede que no llegueos a poder aprovechar este máximo.
    Este parámetro no indica cuál es el máximo que se puede edificar en planta.
    \item Altura máxima.
    La altura máxima, como ya se ha introducido, es otra de las restricciones que impone toda ordenanza si bien, en el caso de la industria (en el que dicha altura máxima suele limitarse bastante -10 a 15 m), normalmente las ordenanzas establecen la posibilidad de excepciones por eigencias del proceso (por ejemplo, una fábrica de ascensores que necesita tener una torre de pruebas, o una torre de craqueo en una refinería).
    Altura máxima es la limitación de altura de la parte más alta de la edificación, que se suele expresas en metros, o en metros y número máximo de plantas de edificio.
\end{itemize}

\subsection{Servicios}
\subsubsection{Gestión de Licencias y Legalizaciones de las Instalaciones}
\paragraph{Actuaciones/Solicitudes a realizar}
\begin{itemize}
    \item Licencia de obras y actividad (Ayuntamiento / Otros organismos)
    \item Autorizaciones particulares previas (Consejerías + otros)
    \item Legalizaciones y Puesta en Servicio (Industria)
    \item Licencia de Primera Ocupación (LPO) (Ayuntamiento previa visita)
    \item Licencia de Funcionamiento (Ayuntamiento previa visita)
    \item Autorizaciones Finales (Consejerías + otros: previa visita)
\end{itemize}

\paragraph{Actos sujetos a Licencia (RD 2187/1978 de 23 junio)}
\begin{itemize}
    \item Obras que hayan de realizarse con carácter provisional.
    \item Obras de instalación de servicios públicos.
    \item Los movimientos de tierra, tales como desmontes, explanación, excavación y terraplenado, salvo que tales actos estén detallados y programados como obtras a ejecutar en un Proyecto de Urbanización o de Edificación aprobado o autorizado.
    \item La primera utilización u ocupación provisional de los edificios e instalaciones en general.
    \item El uso del vuelo sobre las edificaciones e instalaciones de todas clases existentes.
    \item La modificación del uso de los edificios e instalaciones en general.
    \item Demolición de las construcciones, salvo en los casos declarados de ruina inminente.
    \item Las instalaciones subterráneas dedicadas a aparcamientos, actividades industriales, mercantiles o profesionales, servicios públicos o cualquier otro uso a que se destine el subsuelo.
    \item Corta de árboles integrados en masa arbórea que esté enclavada en terrenos para los que exista un Paln de Ordenación aprobado.
    \item Colocación de carteles de propaganda visibles desde la vía pública.
    \item Cuando los actos de edificación y uso del suelo y aquellos otros se realizarán por particulares en terreno de dominio público, sin perjuicio de las autorizaciones o concesiones que sea pertinente otorgar por parte del ente titular del dominio público.
\end{itemize}

\paragraph{Tipos de Licencia}
\begin{itemize}
    \item Licencia de obras actuaciones urbanísticas
    \item Licencia de primera ocupación (cédula de habitabilidad)
    \item Licencia para la instalación de actividades
    \item Licencia de cambio de uso $\rightarrow$ Plan especial
    \item Otras licencia;
    Específicas según el tipo de edificio, autorizaciones sanitarias, industria, terrazas, consejo nuclear, educación, etc \dots
\end{itemize}

\paragraph{Partes interesadas en una concesión de Licencia}
\begin{itemize}
    \item Propietario / Propiedad / Promotor
    \item Dirección Integrada de Proyectos (DIP)
    \item Técnicos habilitados (Arqtos., Ingenieros, etc\dots)
    \item Técnicos de la Administración Pública
    \item Entidades Públicas
    \item Entidades que gestionan por parte de la Admon. las licencias denominadas ECU's (Entidad Colaboradora Urbanística)
    \item Vecinos
    \item Otras entidades o particulares (reclamaciones, denuncias)
    \item OCA: Organismo de Control Autorizado
    \item Industria; Acta de Puesta en servicio.
    \item Registro Industrial, etc... 
\end{itemize}

\paragraph{Procedimientos}
\begin{itemize}
    \item Ordinario: Actuaciones que impliquen obras que afectan a la estructura o a elementos protegidos y por tanto, requieren un proyecto técnico para ser definidas, aprobadas y ejercutadas.
    \item Ordinario común: Actuaciones que implican obras de envergadura, o bien que afecten a elementos protegidos.
    \item Ordinario abreviado: Actuaciones que implican obras de cierta envergadura, pudiendo modificar puntualmente la estructura, pero que en ningún caso afecten a elementos protegidos.
    \item Simplificado: Son aquellas en donde las actividades tienen una limitada influencia en el entorno urbanístico, en el de la seguridad o en su previsible repercusión medioambiental. En este caso una declaración suscrita por técnico competente, visada por el colegio profesional correspondiente, hará constar que la actuación urbanística que se solicita se adecúa a la normativa que le es de aplicación, si bien no será necesario un proyecto técnico.
    \item Actuaciones comunicadas: Aquellas que no conllevan repercusión sobre el entorno, ya sea urbanístico o ambiental; así como aquellas que no requeiren unas condiciones de seguridad especiales. A su vez, las obras necesarias para la implantación de una actividad no afectan a la estructura o a elementos protegidos.
    Por este motivo, no es necesario que un técnico elabore la documentación que tenemos que presentar, aunque en función de la actuación y de nuestras propias capacidades y tiempos puede ser que necesitemos a alguno que nos ayude a elaborar planos o que nos asesore en algún momento.
\end{itemize}

\paragraph{Proceso de Licencias}

\paragraph{Fases de tramitación de Licencias}
\begin{enumerate}
    \item Recabar información necesaria
    \begin{itemize}
        \item Código Técnico de la Edificación (CTE). 
        \item Plan General de Ordenación Urbana del Municipio (PGOU)
        \item Ordenanza de Tramitación de Licencia del Municipio.
        \item Otras ordenanzas aplicables (según el tipo de actividad, etc...)
        \item Ley de Suelo de la Comunidad Autónoma.
        \item Leu de Suelo del Gobierno Estatal
    \end{itemize}
    Si se quiere solicitar cambio de uso hay que realizar un plan especial para aprobación en el Ayuntamiento.
    \item Solicitud de licencia
    \begin{itemize}
        \item Permiso para poder construir lo que contiene el proyecto
        \item Requisitos / Justificación
        \begin{itemize}
            \item Arquitectura:Volumen, edificabilidad, retranqueos, usos permitidos.
            \item Instalaciones: Justificación Medio Ambiente, ruidos, acometida saneamiento, DB-SI.
            \item Gestionar Acometidas: 
            \begin{itemize}
                \item Agua PCI + Agua Fría
                \item Electricidad
                \item comunicaciones
                \item Gas Natural
                \item Otro combustible
            \end{itemize}
        \end{itemize}
    \end{itemize}
    \item Tramitación
    \begin{itemize}
        \item Informe preceptivo (Edificios BIC en Madrid a través de la Comisión Institucional Patrimonio Histórico, Artístico y Natutal (CIPHAN), DB-SI de Protección Civil,movilidad y medio ambiente)
        \item Subsanación de requerimientos (Posible desestimación de Licencia si no se subsanan las deficiencias).
    \end{itemize}
    \item Resolución
    \begin{itemize}
        \item Mediante Certificado de Conformidad emitido por la ECU. (Al cabo de un mes se pueden comenzar las obras)
        \item Concesión final del Ayuntamiento (3 a 6 meses)
    \end{itemize}
    \item Comienzo de las Obras
    \begin{itemize}
        \item Visado del proyecto; es necesaria la Certificación Energética del proyecto para el visado.
        \item Podemos comenzar la construcción.
        \item Tenemos permiso para construir lo aprobado y según la actividad permitida.
        \item Durante el transcurso de la construcción cambios respecto a lo solicitado y aprobado por el Ayuntamiento:
        \begin{itemize}
            \item Solicitud de modificación de licencia de obras y actividad (instalaciones)
            \item Gestionar posibles cambios de normativa
        \end{itemize}
    \end{itemize}
    \item Legalizaciones
    Durante la ejecución de las instalaciones y previo a solicitar la licencia de funcionamiento y dar el CFO de las instalaciones finales se debe solicitar el acta de puesta en marcha en Industria a través de una OCA.
    Instalaciones a legalizar:
    \begin{itemize}
        \item Climatización: Generación de frío o calor $\geq 70 Kw$
        \item Fontanería:
        \begin{itemize}
            \item $\geq$ 25 contadores
            \item Caudal $\geq$ 6 l/seg
            \item $\geq$ 7 fluxores
            \item Acometida $\geq$ 65 mm
            \item $\geq$ 15 plantas
        \end{itemize}
        \item Electricidad:
        \begin{itemize}
            \item Media Tensión
            \item Baja Tensión
            \item Grupos Electrógenos
        \end{itemize}
        \item PCI + Detección de incendios
        \item Fotovoltaicas
        \item Saneamiento
        \item Comunicaciones: ICT
        \item Garajes (instalación BT)
    \end{itemize}
    Se realizan a través de Organismos de Control Autorizados (OCA)
    \begin{itemize}
        \item Pago de las tasas a la OCA
        \item Aprobación por parte de la OCA del proyecto a legalizar
        \item Registro de la documentación en Industria con visado o declaraición responsable
        \item Solicitud visita de la OCA para la aprobación de la ejecución de la instalación
        \item Registro Final en Industria
        \item Obtención del Acta de Puesta en Servicio
    \end{itemize}
    \item Solicitud de la LP y Licencia Funcionamiento
    \begin{itemize}
        \item Firmar CFO
        \begin{itemize}
            \item DF Arquitecto / Estructuras
            \item DF Instalaciones
            \item DE Arquitectura
            \item DE Instalaciones
        \end{itemize}
    \end{itemize}
    \item Autorizaciones finales
    Solicitar permisos:
    \begin{itemize}
        \item Sanidad
        \item Educación
        \item Recogida de Residuos
        \item Otros \dots
    \end{itemize}
\end{enumerate}
